\chapter{Complexity — P Closure}
%%% generated by the blueprint pipeline from TCSlib.Complexity.ClassNP.PClosure
%%%%%%%%%%%%%%%%%%%%%%%%%%%%%%%%%%%%%%%%%%%%%%%%%%%%%%%%%%%%%%%%%%%%%%%%%%%%%%%
\section{Overview}

This module establishes the closure properties of the class $\mathsf{P}$ [AB09, Def.~1.13]
that are used tacitly throughout [AB09]: a language is in $\mathsf{P}$ exactly when its
one-bit indicator is polynomial-time computable, and $\mathsf{P}$ is closed under
polynomial-time preimages (many-one reductions), under intersection and union, and hence
under arbitrary Boolean functions of finitely many tests. As applications, comparing the
lengths of the two components of a pair, or of the outputs of two polynomial-time
functions, is decidable in polynomial time.

%%%%%%%%%%%%%%%%%%%%%%%%%%%%%%%%%%%%%%%%%%%%%%%%%%%%%%%%%%%%%%%%%%%%%%%%%%%%%%%
\section{Declarations}

\begin{theorem}[P via polynomial-time indicator]
\lean{Complexity.mem_P_iff_polyTimeComputable}
\label{Complexity.mem_P_iff_polyTimeComputable}
\leanok
\uses{Complexity.P, Complexity.PolyTimeComputable, Complexity.mem_P_iff, Turing.MultiTapeTM.indicator}
\difficulty{3}
Let $V \subseteq \{0,1\}^*$ be a language and let $\mathbf{1}_V(x) \in \{0,1\}$ be its indicator bit, equal to $1$ if $x \in V$ and $0$ otherwise. Then $V$ belongs to $\mathsf{P}$, the class of languages decidable by a deterministic Turing machine in polynomial time, if and only if the map sending a bit string $x$ to the one-bit word $\mathbf{1}_V(x)$ is polynomial-time computable, i.e.\ some finite binary Turing machine computes it within $C(n+1)^c$ steps on every input of length $n$, for some constants $C, c \in \mathbb{N}$.
\end{theorem}

\begin{theorem}[P is closed under polynomial-time preimages]
\lean{Complexity.preimage_mem_P}
\label{Complexity.preimage_mem_P}
\leanok
\uses{Complexity.P, Complexity.PolyTimeComputable, Complexity.PolyTimeComputable.comp, Complexity.mem_P_iff_polyTimeComputable}
\difficulty{1}
Let $V \subseteq \{0,1\}^*$ be a language in $\mathsf{P}$, the class of languages decidable by a deterministic Turing machine in polynomial time, and let $g : \{0,1\}^* \to \{0,1\}^*$ be polynomial-time computable (computed by some finite binary Turing machine within $C(n+1)^c$ steps on inputs of length $n$, for some constants $C, c$). Then the preimage $g^{-1}(V) = \{x \in \{0,1\}^* : g(x) \in V\}$ also belongs to $\mathsf{P}$.
\end{theorem}

\begin{theorem}[Polynomial-time Boolean tests define languages in P]
\lean{Complexity.mem_P_of_test}
\label{Complexity.mem_P_of_test}
\leanok
\uses{Complexity.P, Complexity.PolyTimeComputable, Complexity.mem_P_iff_polyTimeComputable, Turing.MultiTapeTM.indicator}
\difficulty{3}
Let $b : \{0,1\}^* \to \{0,1\}$ be a Boolean test on bit strings such that the map sending $z$ to the one-bit word $b(z)$ is polynomial-time computable (computed by some finite binary Turing machine within $C(n+1)^c$ steps on inputs of length $n$, for some constants $C, c$). Then the language $\{z \in \{0,1\}^* : b(z) = 1\}$ belongs to $\mathsf{P}$, the class of languages decidable by a deterministic Turing machine in polynomial time.
\end{theorem}

\begin{theorem}[Indicator of a language in P is polynomial-time]
\lean{Complexity.test_of_mem_P}
\label{Complexity.test_of_mem_P}
\leanok
\uses{Complexity.P, Complexity.PolyTimeComputable, Complexity.mem_P_iff_polyTimeComputable, Turing.MultiTapeTM.indicator}
\difficulty{1}
Let $L \subseteq \{0,1\}^*$ be a language in $\mathsf{P}$, the class of languages decidable by a deterministic Turing machine in polynomial time. Then the map sending a bit string $z$ to the one-bit word $\mathbf{1}_L(z)$, which is $1$ if $z \in L$ and $0$ otherwise, is polynomial-time computable, i.e.\ some finite binary Turing machine computes it within $C(n+1)^c$ steps on every input of length $n$, for some constants $C, c \in \mathbb{N}$.
\end{theorem}

\begin{theorem}[P is closed under intersection]
\lean{Complexity.inter_mem_P}
\label{Complexity.inter_mem_P}
\leanok
\uses{Complexity.P, Complexity.mem_P_of_test, Complexity.polyTimeComputable_and, Complexity.test_of_mem_P, Turing.MultiTapeTM.indicator}
\difficulty{3}
Let $L_1, L_2 \subseteq \{0,1\}^*$ be languages in $\mathsf{P}$, the class of languages decidable by a deterministic Turing machine in polynomial time. Then their intersection $L_1 \cap L_2 = \{z : z \in L_1 \text{ and } z \in L_2\}$ also belongs to $\mathsf{P}$.
\end{theorem}

\begin{theorem}[P is closed under union]
\lean{Complexity.union_mem_P}
\label{Complexity.union_mem_P}
\leanok
\uses{Complexity.P, Complexity.PolyTimeComputable, Complexity.mem_P_of_test, Complexity.polyTimeComputable_const, Complexity.polyTimeComputable_ite, Complexity.test_of_mem_P, Turing.MultiTapeTM.indicator}
\difficulty{4}
Let $L_1, L_2 \subseteq \{0,1\}^*$ be languages in $\mathsf{P}$, the class of languages decidable by a deterministic Turing machine in polynomial time. Then their union $L_1 \cup L_2 = \{z : z \in L_1 \text{ or } z \in L_2\}$ also belongs to $\mathsf{P}$.
\end{theorem}

\begin{theorem}[The empty language is in P]
\lean{Complexity.empty_mem_P}
\label{Complexity.empty_mem_P}
\leanok
\uses{Complexity.P, Complexity.mem_P_of_test, Complexity.polyTimeComputable_const}
\difficulty{1}
Let $\mathsf{P}$ be the class of languages over $\{0,1\}$ decidable by a deterministic Turing machine in polynomial time. Then the empty language $\emptyset \subseteq \{0,1\}^*$ belongs to $\mathsf{P}$.
\end{theorem}

\begin{theorem}[The full language is in P]
\lean{Complexity.univ_mem_P}
\label{Complexity.univ_mem_P}
\leanok
\uses{Complexity.P, Complexity.mem_P_of_test, Complexity.polyTimeComputable_const}
\difficulty{1}
Let $\mathsf{P}$ be the class of languages over $\{0,1\}$ decidable by a deterministic Turing machine in polynomial time. Then the language $\{0,1\}^*$ consisting of all bit strings belongs to $\mathsf{P}$.
\end{theorem}

\begin{theorem}[P is closed under Boolean combinations of tests]
\lean{Complexity.mem_P_of_atoms}
\label{Complexity.mem_P_of_atoms}
\leanok
\uses{Complexity.P, Complexity.compl_mem_P, Complexity.empty_mem_P, Complexity.inter_mem_P, Complexity.union_mem_P, Complexity.univ_mem_P}
\difficulty{5}
Let $\iota$ be a finite index set and, for each $i \in \iota$, let $b_i : \{0,1\}^* \to \{0,1\}$ be a Boolean test on bit strings such that the language $\{z : b_i(z) = 1\}$ belongs to $\mathsf{P}$, the class of languages decidable by a deterministic Turing machine in polynomial time. Then for every Boolean function $F : \{0,1\}^{\iota} \to \{0,1\}$, the language
\[ \{\, z \in \{0,1\}^* : F\big((b_i(z))_{i \in \iota}\big) = 1 \,\} \]
also belongs to $\mathsf{P}$.
\end{theorem}

\begin{theorem}[Equal-length pairs form a language in P]
\lean{Complexity.lenEq_mem_P}
\label{Complexity.lenEq_mem_P}
\leanok
\uses{Complexity.P, Complexity.mem_P_of_test, Complexity.pairFstD, Complexity.pairSndD, Complexity.polyTimeComputable_lenEq}
\difficulty{2}
Let $\pi_1$ and $\pi_2$ be the total projections on bit strings, which send a bit string of the form $\langle a,b\rangle$ (the self-delimiting pair encoding: every bit of $a$ doubled, then $01$, then $b$) to $a$ and to $b$ respectively, and send every other bit string to the empty word. Then the language $\{z \in \{0,1\}^* : |\pi_1(z)| = |\pi_2(z)|\}$ belongs to $\mathsf{P}$, the class of languages decidable by a deterministic Turing machine in polynomial time.
\end{theorem}

\begin{theorem}[Length comparison of pair components is in P]
\lean{Complexity.lenLe_mem_P}
\label{Complexity.lenLe_mem_P}
\leanok
\uses{Complexity.P, Complexity.mem_P_of_test, Complexity.pairFstD, Complexity.pairSndD, Complexity.polyTimeComputable_lenLe}
\difficulty{2}
Let $\pi_1$ and $\pi_2$ be the total projections on bit strings, which send a bit string of the form $\langle a,b\rangle$ (the self-delimiting pair encoding: every bit of $a$ doubled, then $01$, then $b$) to $a$ and to $b$ respectively, and send every other bit string to the empty word. Then the language $\{z \in \{0,1\}^* : |\pi_2(z)| \le |\pi_1(z)|\}$ belongs to $\mathsf{P}$, the class of languages decidable by a deterministic Turing machine in polynomial time.
\end{theorem}

\begin{theorem}[Equal output lengths of polynomial-time maps]
\lean{Complexity.lenEq_preimage_mem_P}
\label{Complexity.lenEq_preimage_mem_P}
\leanok
\uses{Complexity.P, Complexity.PolyTimeComputable, Complexity.PolyTimeComputable.pairEncode, Complexity.lenEq_mem_P, Complexity.preimage_mem_P}
\difficulty{2}
Let $f, g : \{0,1\}^* \to \{0,1\}^*$ be polynomial-time computable functions (each computed by some finite binary Turing machine within $C(n+1)^c$ steps on inputs of length $n$, for some constants $C, c$). Then the language $\{z \in \{0,1\}^* : |f(z)| = |g(z)|\}$ belongs to $\mathsf{P}$, the class of languages decidable by a deterministic Turing machine in polynomial time.
\end{theorem}

\begin{theorem}[Output-length comparison of polynomial-time maps]
\lean{Complexity.lenLe_preimage_mem_P}
\label{Complexity.lenLe_preimage_mem_P}
\leanok
\uses{Complexity.P, Complexity.PolyTimeComputable, Complexity.PolyTimeComputable.pairEncode, Complexity.lenLe_mem_P, Complexity.preimage_mem_P}
\difficulty{2}
Let $f, g : \{0,1\}^* \to \{0,1\}^*$ be polynomial-time computable functions (each computed by some finite binary Turing machine within $C(n+1)^c$ steps on inputs of length $n$, for some constants $C, c$). Then the language $\{z \in \{0,1\}^* : |g(z)| \le |f(z)|\}$ belongs to $\mathsf{P}$, the class of languages decidable by a deterministic Turing machine in polynomial time.
\end{theorem}
